\begin{ouverture}[algorithmique]{20}{Penser et programmer par blocs}{Partie 4 · Données, hasard et programmation}
{Entrées · Sorties · Expressions · Séquences · Répétitions simples}
\objectif{Lire et modifier un programme simple.}
\objectif{Représenter une formule par des blocs.}
\objectif{Comprendre une répétition inconditionnelle.}
\prerequis{Formules et angles des polygones rencontrés.}
\end{ouverture}
\pagestyle{otMain}\essentiel\fontsize{11}{15}\selectfont
\phantomsection\label{cours-20-1}
\essSection{1}{Lire une séquence et une formule}
\begin{definition}Un programme est une suite d’instructions exécutées dans un ordre précis. Une entrée est une donnée fournie au programme ; une sortie est une information produite. En Cinquième, une donnée saisie est lue et utilisée sans étudier les mises à jour complexes de variables.\end{definition}
\label{fig-cours-20-1-cours-programme}\figureGeo{cours-programme}

\begin{savoirfaire}[Prévoir avant d’exécuter]Pour une entrée 4, le calcul $3\times\text{réponse}+2$ donne 14. Les blocs imbriqués décrivent les parenthèses : multiplier après addition n’est pas la même chose qu’additionner après multiplication.\end{savoirfaire}
\begin{attention}Le texte \texttt{réponse} représente ici la donnée saisie. Une phrase affichée n’est pas une nouvelle entrée. Identifier où le programme demande, calcule et affiche.\end{attention}
\clearpage
\phantomsection\label{cours-20-2}
\essSection{2}{Répéter une séquence linéaire}
\begin{definition}Une boucle inconditionnelle répète un nombre fixé de fois la même suite d’instructions. « Répéter 4 fois : avancer de 60, tourner de $90^\circ$ » trace un carré si le stylo est baissé.\end{definition}
\label{fig-cours-20-2-cours-boucle}\figureGeo{cours-boucle}

\begin{savoirfaire}[Distinguer angle intérieur et rotation]L’angle du bloc « tourner » est le changement de direction du déplacement. Pour un triangle équilatéral, tourner de $120^\circ$ à chaque sommet, et non de son angle intérieur de $60^\circ$.\end{savoirfaire}
\begin{attention}Vérifier le point de départ, la direction initiale, le stylo et le nombre de répétitions. Une boucle de quatre pas ne trace pas forcément un carré : les longueurs et les angles doivent aussi convenir.\end{attention}

\clearpage\rubrique{Découverte}
\begin{blocLivre}{ex-20-1}\begin{exercice}[1]{Entrée et sortie}
Le programme demande un nombre puis dit « réponse + 5 ». Identifier l’entrée et la sortie pour une saisie de 7.
\label{fig-ex-20-1}\figureGeo{prog-1}
\end{exercice}\end{blocLivre}
\begin{blocLivre}{ex-20-2}\begin{exercice}[1]{Prévoir une formule}
Prévoir la sortie du programme pour une entrée de 4.
\label{fig-ex-20-2}\figureGeo{prog-2}
\end{exercice}\end{blocLivre}
\begin{blocLivre}{ex-20-3}\begin{exercice}[1]{Des parenthèses}
Le programme calcule $3(\text{réponse}+2)$. Quelle sortie donne l’entrée 4 ?
\label{fig-ex-20-3}\figureGeo{prog-3}
\end{exercice}\end{blocLivre}
\begin{blocLivre}{ex-20-4}\begin{exercice}[1]{Une aire}
Le programme calcule le carré de la réponse. Quelle sortie donne une entrée de 5, interprétée en cm ?
\label{fig-ex-20-4}\figureGeo{prog-4}
\end{exercice}\end{blocLivre}
\begin{blocLivre}{ex-20-5}\begin{exercice}[1]{Deux déplacements}
Le lutin, dirigé vers la droite, avance de 30 puis de 50 pas, sans tourner. Quelle distance parcourt-il ?
\label{fig-ex-20-5}\figureGeo{prog-5}
\end{exercice}\end{blocLivre}

\rubrique{Entraînement}
\begin{blocLivre}{ex-20-6}\begin{exercice}[2]{Modifier un paramètre}
Le programme affiche $2\times\text{réponse}+4$. Modifier un seul nombre pour qu’il affiche le triple de la réponse plus 4. Tester avec 6.
\label{fig-ex-20-6}\figureGeo{prog-6}
\end{exercice}\end{blocLivre}
\begin{blocLivre}{ex-20-7}\begin{exercice}[2]{Tracer un carré}
Identifier la figure tracée par « répéter 4 fois : avancer de 50 ; tourner à gauche de $90^\circ$ ». Donner son périmètre en pas.
\label{fig-ex-20-7}\figureGeo{prog-7}
\end{exercice}\end{blocLivre}
\begin{blocLivre}{ex-20-8}\begin{exercice}[2]{Un rectangle en boucle}
Le programme répète deux fois : avancer de 80, tourner de 90 degrés, avancer de 40, tourner de 90 degrés. Identifier la figure et son périmètre.
\label{fig-ex-20-8}\figureGeo{prog-8}
\end{exercice}\end{blocLivre}
\begin{blocLivre}{ex-20-9}\begin{exercice}[2]{Un triangle équilatéral}
Le programme répète trois fois : avancer de 60, tourner de $120^\circ$. Expliquer pourquoi la figure est un triangle équilatéral.
\label{fig-ex-20-9}\figureGeo{prog-9}
\end{exercice}\end{blocLivre}
\begin{blocLivre}{ex-20-10}\begin{exercice}[2]{L’ordre des blocs}
Comparer les trajets « avancer de 40, tourner à gauche de 90 degrés » et « tourner à gauche de 90 degrés, avancer de 40 », avec une direction initiale vers la droite.
\label{fig-ex-20-10}\figureGeo{prog-10}
\end{exercice}\end{blocLivre}

\rubrique{Problèmes}
\begin{blocLivre}{ex-20-11}\begin{exercice}[3]{Un pentagone régulier}
Un script répète cinq fois : avancer de 40, tourner de $72^\circ$. Expliquer la fermeture et calculer le périmètre.
\label{fig-ex-20-11}\figureGeo{prog-11}
\end{exercice}\end{blocLivre}
\begin{blocLivre}{ex-20-12}\begin{exercice}[3]{Une ligne brisée}
Le script répète trois fois : avancer de 30, tourner à gauche de 90 degrés, avancer de 20, tourner à droite de 90 degrés. Décrire le déplacement final.
\label{fig-ex-20-12}\figureGeo{prog-12}
\end{exercice}\end{blocLivre}
\begin{blocLivre}{ex-20-13}\begin{exercice}[3]{Réparer un carré incomplet}
Un script répète trois fois « avancer de 50, tourner de 90 degrés » pour dessiner un carré. Expliquer le défaut et le corriger.
\label{fig-ex-20-13}\figureGeo{prog-13}
\end{exercice}\end{blocLivre}
\begin{blocLivre}{ex-20-14}\begin{exercice}[3]{Une formule de volume}
Un script demande le côté d’un cube et affiche réponse × réponse × réponse. Prévoir les sorties pour 2 et 3, puis préciser les unités si le côté est en cm.
\label{fig-ex-20-14}\figureGeo{prog-14}
\end{exercice}\end{blocLivre}
\begin{blocLivre}{ex-20-15}\begin{exercice}[3]{Concevoir puis contrôler}
Écrire un programme demandant une longueur en cm et affichant le périmètre d’un carré de ce côté. Tester pour 2,5 et 0. Préciser ce que représente le test 0.
\label{fig-ex-20-15}\figureGeo{prog-15}
\end{exercice}\end{blocLivre}

\clearpage
